\chapter{AI 计算在 CPU 上的基本地图}

AI 计算看起来像一串数学公式，落到 CPU 上却是一串非常具体的动作：从内存读参数和激活值，把数据搬进缓存和寄存器，执行乘法、加法、比较、指数、归约，再把结果写回内存。理解这张地图的目的，是避免把“模型很大”“矩阵很多”“代码用了 SIMD”这些模糊说法混在一起。真正能指导优化的，是每一层到底做了多少运算、移动了多少数据、复用了多少次、受哪个硬件资源限制。

\section{从模型公式到机器动作}

以最常见的线性层为例，数学上它写成 $Y = XW + b$。这个公式很短，但 CPU 看到的不是矩阵符号，而是循环、地址、加载、乘加和存储。假设 $X$ 是一批输入，$W$ 是权重矩阵，$Y$ 是输出矩阵，那么每一个输出元素都要遍历一整段输入和权重。朴素代码会把这个过程写成三重循环：外层枚举输出行，中间枚举输出列，内层做点积。

\begin{lstlisting}[language=C++,caption={线性层的朴素三重循环}]
#include <cstddef>
#include <cstdint>
#include <span>

void linear_naive(std::span<const float> x,
                  std::span<const float> w,
                  std::span<const float> bias,
                  std::span<float> y,
                  std::int32_t batch,
                  std::int32_t input_dim,
                  std::int32_t output_dim) {
    for (std::int32_t row = 0; row < batch; ++row) {
        for (std::int32_t col = 0; col < output_dim; ++col) {
            float acc = bias[static_cast<std::size_t>(col)];
            for (std::int32_t k = 0; k < input_dim; ++k) {
                const std::size_t x_index =
                    static_cast<std::size_t>(row) * input_dim + k;
                const std::size_t w_index =
                    static_cast<std::size_t>(k) * output_dim + col;
                acc += x[x_index] * w[w_index];
            }
            y[static_cast<std::size_t>(row) * output_dim + col] = acc;
        }
    }
}
\end{lstlisting}

这段代码的教学价值不在于快，而在于它把问题拆开了。内层循环每走一步，需要读一个输入元素、读一个权重元素、做一次乘法和一次加法。输出元素只有最后写一次，偏置只读一次。如果输入维度很大，那么主要成本来自内层循环。于是第一个问题不是“能不能优化”，而是“内层循环的瓶颈是什么”。如果权重访问跨步很大，缓存命中会差；如果输入被多个输出列重复使用却每次重新加载，复用就没有被充分利用；如果编译器无法向量化，浮点执行单元就吃不满。

\begin{mentalmodel}
AI 算子首先是数据流，其次才是公式。公式告诉我们结果是什么，数据流告诉我们机器为了得到结果必须搬多少数据、保留哪些中间值、在哪些层级复用。
\end{mentalmodel}

\section{操作数、字节数和复用}

高性能计算里常用两个基本量描述一个算子。第一个是操作数，通常用 FLOPs 表示浮点操作数量。一次乘加有时按两个浮点操作计数，因为它包含一次乘法和一次加法；使用 FMA 指令时，硬件可能一条指令完成乘加，但数学工作量仍然按两个操作理解。第二个是字节数，也就是为了完成这些操作至少要从内存层级搬动多少数据。操作数决定计算压力，字节数决定带宽压力。

仅有操作数还不够。一个矩阵乘的 FLOPs 很高，但如果每个权重元素被大量输出复用，内存流量可以被摊薄；一个逐元素激活函数的 FLOPs 不高，但每个元素通常只读一次写一次，复用极少，容易变成内存带宽问题。复用是连接操作数和字节数的关键。相同的公式，数据布局不同，复用位置就不同；相同的布局，循环顺序不同，复用被缓存捕获的概率也不同。

可以用一个非常粗的估算理解线性层。若输入是 $B \times K$，权重是 $K \times N$，输出是 $B \times N$，乘加操作约为 $2BKN$。最理想情况下，输入、权重和输出各读写接近一次，字节数大致与 $(BK + KN + BN)$ 成正比；最糟糕情况下，权重和输入在缓存里留不住，每个输出元素都重新从较慢层级读取，实际流量会接近操作数规模。高性能内核的很多复杂结构，本质上就是让实际流量尽量接近理想下界。

\begin{keyidea}
优化 AI 算子的第一步是写出两个账本：计算账本记录乘法、加法、比较、指数、归约；内存账本记录输入、权重、中间结果和输出的读写字节。没有账本，性能讨论很容易变成猜测。
\end{keyidea}

\section{算子、内核和端到端推理}

讨论 AI 性能时，要区分三个层级。第一层是算子，例如矩阵乘、卷积、Softmax、LayerNorm、Embedding lookup。第二层是内核，也就是某个算子的具体实现，例如 AVX2 的矩阵乘微内核、int8 量化 GEMM、分块 Softmax。第三层是端到端推理，它把多个算子按计算图串起来，还包含内存分配、线程调度、KV Cache 管理、采样和输入输出处理。

一个单独内核很快，不代表端到端推理就快。原因有三类。第一类是形状不匹配，库内核在大矩阵上效率很高，但真实推理可能是 batch 很小、序列长度变化、矩阵很瘦。第二类是中间结果流量，两个算子单独看都快，但中间张量写回内存再读出来，会浪费带宽。第三类是运行时开销，线程池唤醒、任务切分、同步、NUMA 跨节点访问都可能吞掉内核收益。

因此本册不会把“算子优化”和“系统优化”分成互不相干的两件事。我们会先把单个内核讲清楚，再不断问它放进计算图后会发生什么。例如 GEMM 的微内核关注寄存器块，Attention 关注 $QK^T$、Softmax 和 $PV$ 之间的中间矩阵，量化关注每个分组的 scale 和 zero point，运行时关注这些任务怎样切给多个核心。

\section{为什么 CPU 仍然重要}

AI 计算常常和 GPU 绑定在一起，但 CPU 仍然重要。第一，很多部署环境没有独立 GPU，或者功耗、成本、内存容量、启动延迟要求使 CPU 更合适。第二，推理系统总有 CPU 侧工作，例如 tokenization、调度、采样、内存管理、IO 和小算子。第三，理解 CPU 上的性能模型能帮助读者理解更一般的硬件原则：数据搬运昂贵，局部性决定效率，向量化需要规则形状，并行扩展受同步和带宽限制。

CPU 的优势是通用、低延迟、控制流强、内存容量大、系统集成度高；劣势是峰值矩阵吞吐通常低于专用加速器。把 CPU 用好，不是把它想象成小 GPU，而是尊重它的结构：少量强核心、复杂缓存层级、宽 SIMD、强分支预测、成熟操作系统调度和丰富性能计数器。第二册的所有章节都会围绕这些结构展开。

\begin{exercisebox}
选择一个你熟悉的模型层，例如全连接层、卷积层或 Attention 层，写出它的输入、输出和参数形状。先估算 FLOPs，再估算至少需要读写多少字节。不要急着优化，先判断它更像计算受限还是内存受限。
\end{exercisebox}
